%!TEX ROOT = thesis.tex

% =========================================================
% CHAPTER 6: RESULTS
% =========================================================
\chapter{RESULTS}
\label{chap:results}

This chapter reports the empirical results obtained from the revised patient-level validation experiment. Unlike the earlier exploratory scripts, the final experiment used a deterministic 80/20 patient-level split of the ACDC training cohort. The split contained 80 training patients and 20 held-out validation patients, corresponding to 160 training cases and 40 validation cases. The best checkpoint for each model was selected using validation mean Dice, and the final quantitative results were reported only on the held-out validation cases.

The objective of this chapter is descriptive rather than argumentative. It reports what the experiment shows before Chapter 7 interprets the implications. This separation is important because the final results are more nuanced than a simple claim of superiority. Nano-Mamba U-Net is not the highest-Dice model in the final experiment. Its value lies in the combination of competitive segmentation performance and the smallest parameter count among the evaluated models.

\section{Experimental Setting}
All models were trained and evaluated using the same deterministic preprocessing and validation protocol. The input volumes were resized to $256 \times 256 \times 16$, intensity-scaled to $[0,1]$ per volume, and converted to PyTorch tensors using MONAI transforms \cite{cardoso2022monai}. Images used trilinear interpolation, labels used nearest-neighbour interpolation, and no data augmentation or mask post-processing was applied. The evaluated anatomical classes were the Right Ventricle (RV), Myocardium (MYO), and Left Ventricle (LV). The primary segmentation metric was the Dice Similarity Coefficient (DSC), computed on the resized grid.

The comparison included a standard 3D U-Net baseline, Attention U-Net, SegResNet16, the proposed Nano-Mamba U-Net, and two ablation variants. SegResNet16 refers to the configuration with \texttt{init\_filters=16}, which was the version successfully trained under the final patient-level validation protocol. Therefore, the parameter count reported in this chapter is 4.70M for SegResNet16, not the larger 18.80M configuration explored in earlier standalone benchmarking scripts.

For all models, the validation mean Dice was computed as the average of the three foreground classes:
\begin{equation}
    \overline{D}_{val}=\frac{D_{RV}+D_{MYO}+D_{LV}}{3}.
\end{equation}
This formulation avoids allowing one strong class, such as LV, to dominate the interpretation. Reporting RV, MYO, and LV separately is necessary because a clinically useful segmentation model should not only perform well on the easiest structure.

\section{Quantitative Validation Results}
Table \ref{tab:final_results} presents the final held-out validation performance and computational statistics. The results show that SegResNet16 achieved the highest validation mean DSC of 86.70\%. The proposed Nano-Mamba U-Net achieved a mean DSC of 84.78\% with the smallest parameter count among the evaluated models, at 1.46M parameters.

\begin{table}[hbt!]
\caption{Held-out patient-level validation results on the ACDC training cohort. Intervals are post-hoc patient-level percentile-bootstrap 95\% intervals (10,000 replicates); the same validation set was used for checkpoint selection.}
\label{tab:final_results}
\centering
\resizebox{\textwidth}{!}{
\begin{tabular}{l | c c c | c | c | c | c}
\hline
\textbf{Architecture} & \textbf{RV (\%)} & \textbf{MYO (\%)} & \textbf{LV (\%)} & \textbf{Mean DSC (\%)} & \textbf{95\% CI (\%)} & \textbf{Params (M)} & \textbf{FPS}\\
\hline
3D U-Net Baseline & 78.35 & 74.57 & 89.59 & 80.83 & [77.78, 83.46] & 4.81 & 88.17 \\
Attention U-Net & 69.43 & 72.38 & 82.53 & 74.78 & [65.37, 82.44] & 5.91 & 22.41 \\
SegResNet16 & \textbf{84.30} & \textbf{82.76} & \textbf{93.03} & \textbf{86.70} & [84.26, 88.55] & 4.70 & 25.50 \\
\hline
Ablation: No Mamba & 83.26 & 81.30 & 92.37 & 85.64 & [83.08, 87.64] & 2.29 & 28.80 \\
Ablation: Half Mamba & 82.41 & 80.46 & 91.99 & 84.95 & [81.85, 87.21] & 1.64 & 29.04 \\
\hline
Nano-Mamba U-Net & 82.11 & 80.35 & 91.88 & 84.78 & [82.05, 86.90] & \textbf{1.46} & 29.09 \\
\hline
\end{tabular}
}
\end{table}

Compared with the 3D U-Net baseline, Nano-Mamba U-Net improved the mean DSC by 3.95 percentage points while reducing the parameter count from 4.81M to 1.46M. Using the unrounded parameter counts from the result CSV, the reduction relative to 3D U-Net is:
\begin{equation}
    \frac{4.808537-1.456325}{4.808537}\times100\% \approx 69.7\%.
\end{equation}
Compared with SegResNet16, Nano-Mamba U-Net used approximately 69.0\% fewer parameters, but its mean DSC was 1.92 percentage points lower. These results indicate that Nano-Mamba U-Net is not the highest-accuracy model in the final experiment; instead, its main advantage lies in parameter efficiency with competitive validation performance.

Table \ref{tab:paired_bootstrap} reports paired patient-level differences. Positive values favor Nano-Mamba U-Net. Pairing is important because every model was evaluated on the same 20 patients. The intervals are descriptive resampling summaries; no $p$-values, multiplicity correction, or independent test cohort is claimed.

\begin{table}[hbt!]
\caption{Post-hoc paired patient-level Nano-Mamba mean-DSC differences.}
\label{tab:paired_bootstrap}
\centering
\begin{tabular}{l | r | c}
\hline
\textbf{Reference model} & \textbf{Nano minus reference (pp)} & \textbf{95\% bootstrap CI (pp)}\\
\hline
3D U-Net Baseline & +3.95 & [+2.83, +5.01] \\
Attention U-Net & +10.00 & [+3.51, +18.29] \\
SegResNet16 & -1.92 & [-3.01, -0.88] \\
Ablation: No Mamba & -0.86 & [-1.48, -0.23] \\
Ablation: Half Mamba & -0.17 & [-0.91, +0.55] \\
\hline
\end{tabular}
\end{table}

\section{Accuracy-Efficiency Visualization}
Since the research question concerns both segmentation accuracy and model efficiency, the results should not be interpreted from Table \ref{tab:final_results} alone. Figure \ref{fig:accuracy_parameter_tradeoff} visualizes the relationship between validation mean Dice and trainable parameter count. A model positioned higher on the vertical axis has stronger Dice performance, while a model positioned further left has fewer parameters.

\begin{figure}[hbt!]
    \centering
    \includegraphics[width=0.85\textwidth]{accuracy_parameter_tradeoff.png}
    \caption{Accuracy-parameter trade-off generated directly from the audited summary CSV. SegResNet16 obtains the highest held-out validation Dice score, while Nano-Mamba U-Net occupies the lowest-parameter region.}
    \label{fig:accuracy_parameter_tradeoff}
\end{figure}

This visualization makes the central interpretation clearer. SegResNet16 is the accuracy leader and should be acknowledged as such. Nano-Mamba U-Net is not located at the top of the graph, but it appears on the left side because it has the smallest parameter count. The result therefore supports an efficiency-oriented contribution rather than an absolute accuracy claim.

Figure \ref{fig:accuracy_speed_tradeoff} visualizes the relationship between validation mean Dice and inference speed. The speed result should be interpreted carefully because FPS depends on implementation details, hardware, batch size, input resolution, and measurement protocol. Therefore, FPS is reported as a practical engineering indicator rather than as a universal property of the architecture.

\begin{figure}[hbt!]
    \centering
    \includegraphics[width=0.85\textwidth]{accuracy_speed_tradeoff.png}
    \caption{Accuracy-speed trade-off generated directly from the audited summary CSV. The figure provides an engineering view of held-out validation Dice relative to batch-one inference speed in the recorded experiment.}
    \label{fig:accuracy_speed_tradeoff}
\end{figure}

\section{Model-by-Model Interpretation}
The 3D U-Net baseline achieved 80.83\% mean DSC. This result is important because it provides a conventional volumetric segmentation reference point. Nano-Mamba U-Net improved over this baseline on all three structures, suggesting that the proposed compact architecture is not merely smaller but also more accurate than this particular 3D U-Net configuration under the final split.

Attention U-Net achieved the lowest mean DSC among the evaluated models, at 74.78\%. This does not mean attention gates are generally ineffective. It means that, under the implemented settings and validation split, this Attention U-Net configuration did not train or generalize as strongly as the other models. This result also demonstrates why it is unsafe to assume that a more complex architectural concept necessarily improves performance.

SegResNet16 achieved the highest mean DSC, with the best RV, MYO, and LV scores. This suggests that residual convolutional modeling remains very strong for this dataset. Since SegResNet16 outperformed Nano-Mamba U-Net by 1.92 percentage points, the thesis must not describe Nano-Mamba U-Net as the best model by Dice. Instead, SegResNet16 should be acknowledged as the accuracy leader in the final experiment.

The No-Mamba ablation achieved 85.64\% mean DSC, which is 0.86 percentage points higher than Nano-Mamba U-Net. The paired post-hoc interval for Nano minus No-Mamba is $[-1.48,-0.23]$ percentage points. This is an important negative result because it prevents over-claiming the Mamba-inspired module. The No-Mamba design is not parameter matched and uses 2.29M parameters compared with Nano-Mamba's 1.46M, so the comparison does not isolate sequence gating as the cause of the difference. It shows only that the executed No-Mamba architecture obtained higher Dice on this selected split while Nano-Mamba retained the smaller reported parameter count.

The Half-Mamba variant achieved 84.95\% mean DSC, 0.17 percentage points above Nano-Mamba U-Net. The paired interval for Nano minus Half-Mamba is $[-0.91,+0.55]$ percentage points and crosses zero, so the available case-level evidence does not resolve this small ordering reliably. A matched-capacity, multi-seed experiment would be required to study bottleneck size more reliably.

\section{Class-Specific Performance}
Across most models, LV segmentation achieved the highest Dice score, while MYO segmentation was generally more difficult. This pattern is consistent with the anatomical properties of cardiac MRI: the LV cavity is usually more clearly defined, whereas the myocardium is a thinner ring-like structure affected by papillary muscles, partial volume effects, and boundary ambiguity.

For Nano-Mamba U-Net, the LV Dice was 91.88\%, the RV Dice was 82.11\%, and the MYO Dice was 80.35\%. The gap between LV and MYO is:
\begin{equation}
    91.88\%-80.35\%=11.53\%.
\end{equation}
This difference is clinically plausible because the myocardium is more difficult to delineate consistently. A small absolute boundary shift can affect a thin myocardial ring more strongly than a larger ventricular cavity. Therefore, the lower MYO score should be interpreted as a known structural challenge rather than simply as a random model failure.

\begin{figure}[hbt!]
    \centering
    \includegraphics[width=0.90\textwidth]{class_specific_dice_scores.png}
    \caption{Class-specific held-out validation Dice scores generated directly from the audited summary CSV. Nano-Mamba U-Net improves over the evaluated 3D U-Net baseline, while SegResNet16 achieves the strongest scores.}
    \label{fig:dsc_bar_chart}
\end{figure}

\section{Checkpoint Selection Evidence}
The experiment logs contain 150 contiguous epoch rows for every model. For each log, the maximum recorded validation mean DSC occurs at exactly the best epoch in the aggregate table: epochs 105, 121, 134, 149, 62, and 112 for 3D U-Net, Nano-Mamba, No-Mamba, Half-Mamba, Attention U-Net, and SegResNet16, respectively. Figure \ref{fig:training_curves} therefore shows the recorded raw validation traces rather than hand-entered points.

\begin{figure}[hbt!]
    \centering
    \includegraphics[width=0.98\textwidth]{training_validation_curves.png}
    \caption{Raw, unsmoothed 150-epoch validation mean-DSC traces. Markers identify the best validation checkpoint selected for each model. The figure is generated directly from the six experiment-log CSVs.}
    \label{fig:training_curves}
\end{figure}

The learning curves also show why validation-based checkpoint selection was necessary. By epoch 150, the validation score was 0.29 percentage points below the selected maximum for both Nano-Mamba and 3D U-Net, 0.43 points lower for SegResNet16, 0.56--0.62 points lower for the two ablations, and 3.05 points lower for Attention U-Net. Training loss nevertheless continued to decrease for all six models. The larger Attention U-Net gap is the clearest sign of late-epoch overfitting or instability in the recorded runs; it does not imply that attention mechanisms are generally inferior.

The per-case tables identify a shared difficult case. \texttt{patient049\_frame11} is the lowest-mean-Dice case for Nano-Mamba, 3D U-Net, SegResNet16, No-Mamba, and Half-Mamba. Nano-Mamba obtains 57.86\% mean DSC on this case, driven mainly by low RV Dice of 26.04\%. Attention U-Net's worst case is \texttt{patient042\_frame16}, at 16.38\% mean DSC. Because no pathology label or visual overlay is used in this analysis, these numerical failures should not be assigned an anatomical cause without inspecting the corresponding validation images.

No qualitative prediction panel is included in the final thesis. Any future panel should be generated from a validation patient and the selected Nano-Mamba checkpoint, and should clearly state that it illustrates one case rather than proving general performance.

\clearpage

% =========================================================
% CHAPTER 7: DISCUSSION AND LIMITATION
% =========================================================
\chapter{DISCUSSION AND LIMITATION}
\label{chap:discussion}

This chapter interprets the final validation results and explains the implications for the proposed Nano-Mamba U-Net. The discussion is intentionally conservative: the model is evaluated as a lightweight segmentation architecture, not as an absolute state-of-the-art solution. This distinction is necessary because the final results show a trade-off rather than a universal win.

\section{Accuracy-Efficiency Trade-off}
The strongest model by Dice was SegResNet16, with a validation mean DSC of 86.70\%. Nano-Mamba U-Net achieved 84.78\%, which is 1.92 percentage points lower. Therefore, the final experiment does not support the claim that Nano-Mamba U-Net is the most accurate model overall.

However, Nano-Mamba U-Net achieved this performance using only 1.46M parameters, the smallest parameter count among all evaluated models. Relative to the 3D U-Net baseline, Nano-Mamba improved mean DSC by 3.95 percentage points while using substantially fewer parameters. This supports the central contribution of the thesis: Nano-Mamba U-Net provides a compact architecture with competitive validation performance, suitable for exploring lightweight cardiac segmentation.

The trade-off can be interpreted through the relationship between model size and segmentation accuracy. If $D$ denotes mean Dice and $P$ denotes parameter count, then an informal efficiency ratio can be written as:
\begin{equation}
    E = \frac{D}{P}.
\end{equation}
This ratio should not be treated as a clinical metric, but it helps illustrate the engineering perspective. Nano-Mamba U-Net has a lower Dice than SegResNet16 but also a much lower parameter count. Therefore, its value is not that it dominates every metric, but that it occupies a favorable lightweight region of the accuracy-efficiency space.

\section{Interpretation of Ablation Results}
The ablation results require careful interpretation. The No-Mamba variant achieved 85.64\% mean DSC, slightly higher than Nano-Mamba U-Net. This indicates that the current lightweight Mamba-inspired bottleneck does not consistently improve Dice over a stronger convolutional bottleneck under the final validation split. Therefore, the Mamba-inspired module should not be presented as conclusively superior in segmentation accuracy.

At the same time, Nano-Mamba U-Net reduced the parameter count compared with the No-Mamba variant, from 2.29M to 1.46M, while maintaining a close mean DSC. This suggests that the bottleneck is more valuable as a parameter-efficiency design than as a guaranteed accuracy-improving component. The Half-Mamba variant achieved 84.95\%, which is close to Nano-Mamba, further showing that channel allocation and bottleneck capacity affect the accuracy-efficiency balance.

A possible explanation is that the compressed bottleneck operates on a very low-resolution representation. At this stage, the network may already have lost some fine anatomical boundary information. A sequence module can mix information across the compressed representation, but it cannot recover details that were not preserved by the encoder. This explains why the Mamba-inspired module may help compactness without guaranteeing the best Dice. Future work should investigate whether placing sequence modules at multiple resolutions would improve this balance.

\section{Class-Specific Observations}
LV segmentation was the easiest class across the main models, while MYO segmentation was more difficult. This pattern is clinically reasonable because MYO is a thinner structure and its boundary can be affected by papillary muscles and partial volume effects. RV segmentation also remains challenging because the right ventricle has a more variable shape than the left ventricle, especially in pathological cases.

The class-specific findings support the need for reporting RV, MYO, and LV separately. If only mean Dice were reported, the relatively strong LV score could hide weaker MYO performance. Since cardiac measurements such as myocardial mass depend on MYO segmentation, MYO accuracy cannot be treated as a secondary detail. A future version of the project should therefore include additional metrics such as Hausdorff Distance or Average Surface Distance to better evaluate boundary quality.

\section{Relevance to Cardiac Motion Analysis}
The thesis title includes cardiac motion analysis, but the implemented experiment is segmentation-based. The model does not estimate a dense displacement field, optical flow, temporal deformation map, or temporal consistency. Instead, it segments labelled ED and ES 3D volumes as separate cases. In the input tensor, $D$ denotes anatomical slice depth rather than time. The resulting masks could support later computation of phase-specific volumes, but this experiment does not calculate ejection fraction or other motion-related clinical indices.

This limitation does not invalidate the project, but it narrows the claim. The correct interpretation is that Nano-Mamba U-Net contributes to segmentation as a prerequisite for cardiac functional analysis. A stronger future version would process complete 4D cine-MRI sequences and evaluate temporal consistency directly.

\section{Limitations}
Several limitations must be acknowledged.

\begin{enumerate}
    \item \textbf{Internal validation and checkpoint reuse:} The final experiment used a deterministic patient-level validation split of the ACDC training cohort. The same validation patients were inspected across epochs to select each checkpoint and were then used for the descriptive result table. This is not an independent test set and does not represent official ACDC challenge performance.
    \item \textbf{Single split, seed, and post-hoc uncertainty:} Only seed 42 and one 80/20 patient split were reported. The available per-case CSVs support patient-level bootstrap intervals and paired differences, but these were computed after the experiment and use the same validation patients involved in checkpoint selection. They do not replace multi-seed training or an independent test cohort.
    \item \textbf{Mamba-inspired implementation:} The implemented bottleneck is inspired by Mamba-style sequence modeling but is not a full reproduction of the original hardware-aware selective scan implementation. This distinction is important because the thesis must not attribute the full original Mamba mechanism to the implemented model.
    \item \textbf{Legacy projection parameters:} The historical gate projection produces $2d_{state}+1$ channels but uses only the first. The reported parameter count includes the unused outputs. They were retained to preserve checkpoint compatibility, but a redesigned block would need retraining.
    \item \textbf{No explicit temporal modeling:} The implemented task segments ED and ES 3D MRI volumes independently. It performs neither dense motion tracking nor sequence-level consistency analysis, and the depth dimension is not time.
    \item \textbf{Non-matched training and ablations:} Memory constraints required batch size 1 for Attention U-Net and SegResNet16 and batch size 2 for the other models. Optimizer updates per epoch and normalization choices therefore differ across model families. The No-Mamba ablation also has more parameters and a structurally different bottleneck. These factors limit causal attribution of performance differences to one architectural mechanism.
    \item \textbf{Array-size preprocessing and no augmentation:} The pipeline resizes arrays to $256\times256\times16$ but does not explicitly canonicalize orientation, preserve a common physical voxel spacing, crop around the heart, or inverse-resample predictions to native geometry. No training-time data augmentation is used. Dice is therefore measured on a deterministic resized grid, and geometric robustness is not tested.
    \item \textbf{Metric convention:} An empty prediction and empty reference for one class are assigned Dice 1 in the executed metric. None of the 720 audited foreground class scores equals 1, establishing that this branch did not affect the reported case table; the convention could still matter in future datasets.
    \item \textbf{No phase- or pathology-stratified analysis:} Although ACDC contains ED/ES phase information and pathology labels, the final experiment pools the two cases per validation patient and reports aggregate performance. It does not establish whether performance differs systematically between ED and ES or among diagnostic groups.
    \item \textbf{Limited external validation:} The model was not evaluated on a separate multi-center external dataset. Generalization to other scanners and institutions remains future work.
    \item \textbf{Limited metrics and qualitative evidence:} The main reported metric is Dice; no Hausdorff Distance, Average Surface Distance, calibration metric, or blinded qualitative review is available. FPS remains a hardware- and implementation-specific engineering indicator rather than a universal benchmark.
\end{enumerate}

These limitations define the thesis boundary. The case rows and training logs are internally consistent with the numerical table and post-hoc intervals, but they do not remove the design limitations above. Clinical deployment, temporal motion analysis, and external generalization remain outside scope.

\section{Summary of Findings}
The final validation results reposition Nano-Mamba U-Net as a lightweight, parameter-efficient segmentation model. It outperformed the 3D U-Net baseline in mean DSC and used fewer parameters than all evaluated alternatives. Nevertheless, SegResNet16 achieved the highest absolute Dice score, and the No-Mamba ablation slightly exceeded Nano-Mamba in mean DSC. These findings support a balanced conclusion: the proposed model is competitive and efficient, but not definitively superior in accuracy.
